\documentclass[../../main.tex]{subfiles}
\begin{document}
\section{流形上的微分形式}\label{section:ch14-differential-forms-manifolds}

现在我们把注意力转向一个 $n$ 维光滑流形 $M$（带边或不带边）。回忆一下，$\otimes^k T^*M$ 是 $M$ 上协变 $k$-张量丛。其中由斜对称张量组成的子集记作 $\Lambda^k T^*M$：
\begin{align*}
\Lambda^k T^*M=\bigcup_{p\in M}\Lambda^k(T^*_pM).
\end{align*}

\begin{exercise}\label{exercise:ch14-14}
证明 $\Lambda^k T^*M$ 是 $\otimes^k T^*M$ 的一个光滑子丛，因而是 $M$ 上一个秩为 $\binom{n}{k}$ 的光滑向量丛。
\end{exercise}

$\Lambda^k T^*M$ 的一个截面称为一个\emph{微分 $k$-形式}，或简称 $k$-形式；它是一个（连续的）张量场，其在每一点的值都是一个斜对称张量。整数 $k$ 称为形式的\emph{次数}。我们把光滑 $k$-形式所成的向量空间记作
\begin{align*}
\Omega^k(M)=\Gamma\bigl(\Lambda^k T^*M\bigr).
\end{align*}
两个微分形式的外积逐点定义：$(\omega\wedge\eta)_p=\omega_p\wedge\eta_p$。因此，一个 $k$-形式与一个 $l$-形式的外积是一个 $(k+l)$-形式。若 $f$ 是一个 $0$-形式、$\eta$ 是一个 $k$-形式，则我们把外积 $f\wedge\eta$ 理解为通常的数乘 $f\eta$。若定义
\begin{align}
\Omega^*(M)=\bigoplus_{k=0}^n\Omega^k(M),\label{align:ch14-omega-algebra}
\end{align}
则外积使 $\Omega^*(M)$ 成为一个结合的、反交换的分次代数。

在任意光滑坐标图中，一个 $k$-形式 $\omega$ 可局部写成
\begin{align*}
\omega=\sideset{}{'}\sum_I\omega_I\,\mathrm{d}x^{i_1}\wedge\cdots\wedge\mathrm{d}x^{i_k}
=\sideset{}{'}\sum_I\omega_I\,\mathrm{d}x^I,
\end{align*}
其中系数 $\omega_I$ 是定义在坐标域上的连续函数，而我们用 $\mathrm{d}x^I$ 作为 $\mathrm{d}x^{i_1}\wedge\cdots\wedge\mathrm{d}x^{i_k}$ 的缩写（不要与实值函数 $x^I$ 的微分混淆）。命题\ref{proposition:ch12-smooth-components}表明，$\omega$ 在 $U$ 上光滑当且仅当分量函数 $\omega_I$ 都光滑。用微分形式的话说，引理\ref{lemma:ch14-elementary-props}(c)的结果可译为
\begin{align*}
\mathrm{d}x^{i_1}\wedge\cdots\wedge\mathrm{d}x^{i_k}
\left(\frac{\partial}{\partial x^{j_1}},\dots,\frac{\partial}{\partial x^{j_k}}\right)
=\delta^I_J.
\end{align*}
因此 $\omega$ 的分量函数 $\omega_I$ 由
\begin{align*}
\omega_I=\omega\left(\frac{\partial}{\partial x^{i_1}},\dots,\frac{\partial}{\partial x^{i_k}}\right)
\end{align*}
确定。

\begin{example}\label{example:ch14-15forms}
一个 $0$-形式就是一个连续实值函数，一个 $1$-形式就是一个余向量场。在 $\mathbb{R}^3$ 上，光滑 $2$-形式的几个例子是
\begin{align*}
\omega=(\sin xy)\,\mathrm{d}y\wedge\mathrm{d}z;\qquad
\eta=\mathrm{d}x\wedge\mathrm{d}y+\mathrm{d}x\wedge\mathrm{d}z+\mathrm{d}y\wedge\mathrm{d}z.
\end{align*}
$\mathbb{R}^3$ 上每个 $3$-形式都是某个连续实值函数乘以 $\mathrm{d}x\wedge\mathrm{d}y\wedge\mathrm{d}z$。
\end{example}

若 $F:M\to N$ 是一个光滑映射，$\omega$ 是 $N$ 上的一个微分形式，则拉回 $F^*\omega$ 是 $M$ 上的一个微分形式；其定义与任意协变张量场相同：
\begin{align*}
(F^*\omega)_p(v_1,\dots,v_k)=\omega_{F(p)}\bigl(\mathrm{d}F_p(v_1),\dots,\mathrm{d}F_p(v_k)\bigr).
\end{align*}

\begin{lemma}\label{lemma:ch14-pullback-forms}
设 $F:M\to N$ 光滑。
\begin{enumerate}[(a)]
\item $F^*:\Omega^k(N)\to\Omega^k(M)$ 是 $\mathbb{R}$-线性的。
\item $F^*(\omega\wedge\eta)=(F^*\omega)\wedge(F^*\eta)$。
\item 在任意光滑坐标图中，
\begin{align*}
F^*\left(\sideset{}{'}\sum_I\omega_I\,\mathrm{d}y^{i_1}\wedge\cdots\wedge\mathrm{d}y^{i_k}\right)
=\sideset{}{'}\sum_I(\omega_I\circ F)\,
\mathrm{d}(y^{i_1}\circ F)\wedge\cdots\wedge\mathrm{d}(y^{i_k}\circ F).
\end{align*}
\end{enumerate}
\end{lemma}

\begin{exercise}\label{exercise:ch14-17}
证明上一引理。
\end{exercise}

这个引理给出了微分形式拉回的一条计算法则，与我们早先为余向量场和任意张量场所建立的法则类似。

\begin{example}\label{example:ch14-18pullback}
由 $F(u,v)=(u,v,u^2-v^2)$ 定义 $F:\mathbb{R}^2\to\mathbb{R}^3$，令 $\omega$ 为 $\mathbb{R}^3$ 上的 $2$-形式 $y\,\mathrm{d}x\wedge\mathrm{d}z+x\,\mathrm{d}y\wedge\mathrm{d}z$。拉回 $F^*\omega$ 计算如下：
\begin{align*}
F^*(y\,\mathrm{d}x\wedge\mathrm{d}z+x\,\mathrm{d}y\wedge\mathrm{d}z)
&=v\,\mathrm{d}u\wedge\mathrm{d}(u^2-v^2)+u\,\mathrm{d}v\wedge\mathrm{d}(u^2-v^2)\\
&=v\,\mathrm{d}u\wedge(2u\,\mathrm{d}u-2v\,\mathrm{d}v)+u\,\mathrm{d}v\wedge(2u\,\mathrm{d}u-2v\,\mathrm{d}v)\\
&=-2v^2\,\mathrm{d}u\wedge\mathrm{d}v+2u^2\,\mathrm{d}v\wedge\mathrm{d}u,
\end{align*}
其中我们用了反交换性所给出的事实 $\mathrm{d}u\wedge\mathrm{d}u=\mathrm{d}v\wedge\mathrm{d}v=0$。由于 $\mathrm{d}v\wedge\mathrm{d}u=-\mathrm{d}u\wedge\mathrm{d}v$，上式化简为
\begin{align*}
F^*\omega=-2(u^2+v^2)\,\mathrm{d}u\wedge\mathrm{d}v.
\end{align*}
\end{example}

同样的技巧也可用于计算一个微分形式在另一个光滑坐标图中的表达式。

\begin{example}\label{example:ch14-19polar}
令 $\omega=\mathrm{d}x\wedge\mathrm{d}y$ 为 $\mathbb{R}^2$ 上的形式。把到极坐标的变换 $x=r\cos\theta$、$y=r\sin\theta$ 视为恒等映射在定义域与陪域上分别采用不同坐标的表达式，便得
\begin{align*}
\mathrm{d}x\wedge\mathrm{d}y
&=\mathrm{d}(r\cos\theta)\wedge\mathrm{d}(r\sin\theta)\\
&=(\cos\theta\,\mathrm{d}r-r\sin\theta\,\mathrm{d}\theta)\wedge(\sin\theta\,\mathrm{d}r+r\cos\theta\,\mathrm{d}\theta)\\
&=r\,\mathrm{d}r\wedge\mathrm{d}\theta.
\end{align*}
\end{example}

这个公式与把二重积分从直角坐标变换到极坐标的公式之间的相似性是惊人的。下面的命题把这一点推广开来。

\begin{proposition}[顶次形式的拉回公式]\label{proposition:ch14-pullback-top-degree}
设 $F:M\to N$ 是带边或不带边的 $n$ 流形之间的一个光滑映射。若 $(x^i)$ 与 $(y^j)$ 分别是开子集 $U\subseteq M$、$V\subseteq N$ 上的光滑坐标，$u$ 是 $V$ 上的连续实值函数，则在 $U\cap F^{-1}(V)$ 上有
\begin{align}
F^*(u\,\mathrm{d}y^1\wedge\cdots\wedge\mathrm{d}y^n)
=(u\circ F)(\det DF)\,\mathrm{d}x^1\wedge\cdots\wedge\mathrm{d}x^n,\label{align:ch14-top-pullback}
\end{align}
其中 $DF$ 表示 $F$ 在这些坐标下的 Jacobian 矩阵。
\end{proposition}

\begin{proof}
由于 $\Lambda^n T^*M$ 的纤维在每一点都由 $\mathrm{d}x^1\wedge\cdots\wedge\mathrm{d}x^n$ 张成，只需证明把\eqref{align:ch14-top-pullback}两边作用于 $(\partial/\partial x^1,\dots,\partial/\partial x^n)$ 时给出相同结果即可。由引理\ref{lemma:ch14-pullback-forms}，
\begin{align*}
F^*(u\,\mathrm{d}y^1\wedge\cdots\wedge\mathrm{d}y^n)
=(u\circ F)\,\mathrm{d}F^1\wedge\cdots\wedge\mathrm{d}F^n.
\end{align*}
命题\ref{proposition:ch14-wedge-props}(e)表明
\begin{align*}
\mathrm{d}F^1\wedge\cdots\wedge\mathrm{d}F^n\left(\frac{\partial}{\partial x^1},\dots,\frac{\partial}{\partial x^n}\right)
=\det\left(\mathrm{d}F^j\left(\frac{\partial}{\partial x^i}\right)\right)
=\det\left(\frac{\partial F^j}{\partial x^i}\right).
\end{align*}
因此，\eqref{align:ch14-top-pullback}的左端在作用于 $(\partial/\partial x^1,\dots,\partial/\partial x^n)$ 时给出 $(u\circ F)\det DF$。另一方面，右端给出同样结果，因为 $\mathrm{d}x^1\wedge\cdots\wedge\mathrm{d}x^n(\partial/\partial x^1,\dots,\partial/\partial x^n)=1$。
\end{proof}

\begin{corollary}\label{corollary:ch14-top-change-coords}
若 $(U,(x^i))$ 与 $(\widetilde U,(\widetilde x^j))$ 是 $M$ 上两个相交的光滑坐标图，则在 $U\cap\widetilde U$ 上有
\begin{align}
\mathrm{d}\widetilde x^1\wedge\cdots\wedge\mathrm{d}\widetilde x^n
=\det\left(\frac{\partial\widetilde x^j}{\partial x^i}\right)\,
\mathrm{d}x^1\wedge\cdots\wedge\mathrm{d}x^n.\label{align:ch14-top-change}
\end{align}
\end{corollary}

\begin{proof}
把上一命题用于等于 $U\cap\widetilde U$ 上恒等映射的 $F$，但在定义域用坐标 $(x^i)$、在陪域用坐标 $(\widetilde x^j)$。
\end{proof}

内乘也自然地推广到向量场与微分形式，只需逐点作用即可：若 $X\in\mathfrak X(M)$、$\omega\in\Omega^k(M)$，定义一个 $(k-1)$-形式 $X\lrcorner\omega=i_X\omega$：
\begin{align*}
(X\lrcorner\omega)_p=X_p\lrcorner\omega_p.
\end{align*}

\begin{exercise}\label{exercise:ch14-22}
设 $X$ 是 $M$ 上的一个光滑向量场。
\begin{enumerate}[(a)]
\item 证明若 $\omega$ 是一个光滑微分形式，则 $i_X\omega$ 光滑。
\item 验证 $i_X:\Omega^k(M)\to\Omega^{k-1}(M)$ 关于 $C^\infty(M)$ 是线性的，因而对应于一个光滑丛同态 $i_X:\Lambda^k T^*M\to\Lambda^{k-1}T^*M$。
\end{enumerate}
\end{exercise}
\end{document}
